This study elucidates the interfacial dynamics within Ti--Al systems by combining MD simulations with MSM analyses. Our MD simulations reveal that under standard heat treatment temperatures, the interfacial Al liquefies and diffuses through TiAl\textsubscript{3} grain boundaries to the Ti/TiAl\textsubscript{3} interface, while the diffusivity of interfacial Ti remains considerably low. This results in dense accumulation of Al near the Ti/TiAl\textsubscript{3} interface, indicating that the predominant process during heat treatment is the migration of Al atoms towards this interface. This rapid migration suggests a faster formation of TiAl\textsubscript{3} at the Ti/TiAl\textsubscript{3} interface as compared to the Al/TiAl\textsubscript{3} interface. Through MSM analyses, we identified three distinct dynamic states of Al atoms: (1) Al atoms proximate to the substrate, (2) Al atoms at the interfacial region, and (3) Al atoms within the liquid mixture. By examining the eigenvectors and timescales of the dynamic processes, the slowest transition was between state 1 and 3, with the next slowest occurring between state 2 and 3. Furthermore, our findings suggest that with adequate time, as Ti diffuses gradually into the mixture, the resulting increase in Ti concentration triggers a rapid deceleration of all dynamic processes. The dynamic processes typically exhibit slower rates near the Ti substrate compared to the Al substrate. Additionally, a marginal reduction in heat treatment temperature can significantly retard all the dynamic processes. Taken together, the MD and MSM analyses provide insights into the diffusion mechanism in the Ti--Al system, characterized as a three-stage process. Initially, Al atoms diffuse into the TiAl\textsubscript{3} grain boundaries due to a premelting effect, thereby facilitating Al's migration towards the Ti phase. Subsequently, Ti atoms diffuse into the liquid Al, resulting in a mixture where Al atoms predominantly deposit onto the Ti surface. Eventually, this process terminates when the concentration of Ti in the mixture reaches a specific limit. With these insights into the diffusion mechanisms, we hope to inspire new ways to improve the control and optimization of manufacturing processes for these high-performance Ti--Al based materials to minimize defects or enhance their structural properties.
% \jj{With these insights into the diffusion mechanisms, we hope to inspire new ways to improve the control and optimization of manufacturing processes for these high-performance Ti-Al based materials to minimize defects or enhance their structural properties.}